 % Listings (programming language)


\makeatletter
\let\orig@lstnumber=\thelstnumber
\newcommand\lstsetnumber[1]{\gdef\thelstnumber{#1}}
\newcommand\lstresetnumber{\global\let\thelstnumber=\orig@lstnumber}
\makeatother

\lstset{
%   float=hbt,
%   lineskip=-1pt,
%   xrightmargin=1cm,
%   aboveskip=1ex,
%   belowskip=1ex,
    basewidth={0.5em,0.45em}, % reduce flexible column width for characters
    %backgroundcolor=\color{gray!10}, % background fill color
    basicstyle=\ttfamily, % font style used for code
    breakatwhitespace=true, % sets if automatic breaks should only happen at whitespace
    breaklines=true, % sets automatic line breaking
    %caption=\lstname, % set listing caption to filename if \lstinputlisting{filename} is used
    %captionpos=b, % sets the caption position to bottom
    commentstyle=\color{gray!75!black}, % comment style
	escapechar=@,
    frame=single, % frame style 
    framerule=0.05em, % frame rule thickness
    framextopmargin=2pt, % add some margin to frame on top
    framexbottommargin=2pt, % add some margin to frame on bottom
	keywordstyle=[1]\bfseries\color{blue}, % keyword style
	keywordstyle=[2]\bfseries\color{green!50!black},
	keywordstyle=[3]\bfseries\color{red!50!black},
	keywordstyle=[4]\bfseries\color{magenta},
    numbers=left, % where to put the line-numbers; possible values: none, left, right
    numberstyle=\sffamily\footnotesize\color{gray!75!black}, % font style used for line numbers
    rulecolor=\color{gray!55}, % frame rule color
    showstringspaces=false, % mark spaces in strings
    stringstyle=\color{red}, % string literal style
    tabsize=2, % convert any tab into such many whitespaces
    xleftmargin=2em, % makes line numbers appear inside the text block
    %% correct display of German special characters, if needed:
    inputencoding=utf8,
    extendedchars=true,
    %literate={ß}{{\ss}}1 {Ä}{{\"A}}1 {Ö}{{\"O}}1 {Ü}{{\"U}}1 {ä}{{\"a}}1 {ö}{{\"o}}1 {ü}{{\"u}}1,
}
\lstdefinelanguage{CiviC}{
	morekeywords=[1]{extern, export, void, int, for, return},
	morekeywords=[2]{printInt, scanInt, main, readValues, printValues},
	morekeywords=[3]{i, n, v, values}
}
\lstdefinelanguage{CivAS}{
	alsoletter={.,:,_,0,1,2,3,4,5,6,7,8,9},
	morekeywords=[1]{
		.exportfun, .importfun,
		esr,
		isrg,
		jsre,
		jump,
		branch_f,
		iinc_1,
		istore,
		istorea,
		iloada,
		iload_0,
		iload_1,
		iload_2,
		iload_3,
		inewa,
		astore,
		aload_1,
		aload_2,
		jsr,
		iloadc_0,
		return,
		ireturn,
		ilt},
	morekeywords=[2]{
		main:, main,
		readValues:, readValues,
		3_while:, 3_while,
		4_end:, 4_end,
		printValues:, printValues,
		1_while:, 1_while,
		2_end:, 2_end
	},
	otherkeywords={->, \&},
	morekeywords=[3]{int, void},
	morekeywords=[4]{0,1,2,3,4,5,6,7,8,9},
	morestring=[b]"
}
